
\section{Simony}

Editorial article in "Foldebladet" for May 18, 1898.

Since the days of Simon Magus (Acts 8), one has called it simony when someone has used money and outward means to obtain that which can rightly be obtained only by spiritual means. To buy an office in the Church is simony; to buy and sell congregations is doubly so.

This latter form of simony has become rather widespread by the help of the ties of the church bodies and the party zeal that follows with them. In some places it is carried on in a wholly coarse and open manner. There is a small congregation in a place that in church-political significance is called important. The congregation has been formed according to the old church order; one has asked only who wants a pastor, and in no way asked whether one wants a congregation. Naturally the congregation is beforehand well fitted for the work of church politics. If then occasion arises for the choosing of a pastor, agitation breaks loose from within and from without.

The matter is this: from which church body shall one call a pastor? And if the congregation thinks it is too poor to support a pastor by itself, then the matter is where one shall turn for help from the mission fund. Here lies the opportunity of the church parties. They have money, gathered for the purpose of working for the spread of the Kingdom of God; the question of where this money is to be used is often difficult to decide. If, then, the power to decide this lies in the hands of the great church politicians, one knows how it will go.

A congregational meeting is held for the choosing of a pastor in such a small and poor congregation. When the deliberations have begun, one of the politicians in the congregation comes forward with a letter from one of the politicians in one of the church bodies, stating that since the congregation is poor, but loyal to the church body and has bravely stood against "fanaticism," the church body will help the congregation by paying its pastor from the mission fund, provided that the mission committee is allowed to decide which pastor the congregation shall have.

That is a telling stroke in the matter, and no small one either. For a congregation in which a majority is perhaps altogether worldly is very susceptible to worldly arguments. If one was in doubt before, the majority is convinced now, and with thanks one accepts the offer and turns to the kindly disposed mission committee. The congregation has been won. — Simony has been committed.

Such practice is carried on among us, more and less crudely, and thereby irreparable damage is done to congregational life. In this way congregations surrender their most sacred rights and their Christian independence, at the same time that the mission money becomes a means of corruption. It is one of the most important tasks of the Free Church to awaken the right understanding and sense for a Christian work of home mission.

Also on this point there is only one solution that is Christian. It is that the work of home mission direct its aim toward congregation, not toward church body.* Only by self-denying and self-sacrificing labor for the building up of the congregation in each place, without regard to what the "party" may gain, does one truly do any good.

Georg Sverdrup
